\chapter{Estimation: Method of Moments and Maximum Likelihood}

So far this book has run in one direction. You knew the coin was fair, the
machine emitted photons at a rate of three per second, the errors were Gaussian
with mean zero and standard deviation two -- and from those known truths you
deduced what the data would look like. Probability flows from model to data.

Real life runs the other way. You are handed the data first. A new coin lands
heads $61$ times in $100$ tosses. A call center logs the gaps between incoming
calls. A drug trial reports recovery in $73$ of $120$ patients. Now you must
swim upstream: from the data, infer the model. What is the coin's true bias?
What is the call rate? Does the drug help, and by how much? This upstream
journey is \emph{statistical estimation}, and it is where probability stops
being a game about dice and starts being the engine of empirical science.

This chapter introduces the two oldest and most important estimation engines:
the \emph{method of moments}, which matches the averages you see to the
averages the model predicts, and \emph{maximum likelihood}, which asks a
disarmingly simple question, ``which parameter value would have made my data
most plausible?'' We then confront the next questions. An estimate
is itself random -- it dances around as the data dances. Is it centered on the
truth (bias)? How wildly does it dance (variance)? Does it home in on the
truth as data piles up (consistency)? And is there a floor below which no
estimator's dance can shrink (the Cramer--Rao bound)? By the end you should not
just know the formulas; you should see estimation as a single coherent story
about extracting a signal from noise.

\chapterline

\section{Statistical Models and Parameters}
\label{ch30:models}

Before you can estimate anything you must say what you are estimating. That
requires a \emph{statistical model}: a family of probability distributions, one
of which we believe (or pretend) generated the data.

\begin{definition}[Statistical model and parameter]
\label{ch30:def-model}
A \emph{parametric statistical model} is a family of distributions
$\set{P_\theta : \theta \in \Theta}$ indexed by a \emph{parameter} $\theta$
living in a \emph{parameter space} $\Theta \subseteq \R^k$. We assume the
observed data $X_1,\dots,X_n$ are drawn from $P_{\theta_0}$ for some
\emph{true but unknown} $\theta_0 \in \Theta$. An \emph{estimator} is any
function $\hat\theta = g(X_1,\dots,X_n)$ of the data used to guess $\theta_0$.
\end{definition}

Read that definition slowly, because three distinct objects are hiding in it,
and confusing them is the single most common source of muddle in all of
statistics.

\begin{itemize}
\item The \emph{parameter} $\theta_0$ is a fixed, non-random number. The coin
  really does have some bias; we just do not know it. It does not have a
  distribution. It is what it is.
\item The \emph{data} $X_1,\dots,X_n$ are random -- before you collect them they
  could come out many ways.
\item The \emph{estimator} $\hat\theta$ is a function of the random data, so it
  too is random. It is a random variable. It has a distribution, a mean, a
  variance. That distribution -- called the \emph{sampling distribution} of the
  estimator -- is the central character of this chapter.
\end{itemize}

\begin{focusbox}
A parameter is a fixed unknown number. An estimator is a random variable that
tries to pin it down. Almost everything that follows is the study of \emph{how
well a random thing can pin down a fixed thing}.
\end{focusbox}

A small notational convention, used throughout statistics and throughout this
book: a hat means ``estimate of.'' So $\hat\theta$ estimates $\theta$,
$\hat p$ estimates $p$, $\hat\sigma^2$ estimates $\sigma^2$. When we want to
emphasize the estimator depends on $n$ data points we may write $\hat\theta_n$.

\begin{example}[Three running models]
\label{ch30:ex-three-models}
Keep these three in mind; we return to each repeatedly.
\begin{enumerate}
\item \textbf{Bernoulli / coin.} $X_i \sim \Bern(p)$, so $X_i \in \set{0,1}$
  with $\P(X_i=1)=p$. Here $\theta = p \in [0,1]$. Data: $61$ heads in $100$.
\item \textbf{Poisson / counts.} $X_i \sim \Pois(\lambda)$, the number of calls
  in successive minutes. Here $\theta = \lambda > 0$.
\item \textbf{Normal.} $X_i \sim \Normal(\mu,\sigma^2)$. Here
  $\theta = (\mu,\sigma^2)$ is two-dimensional, $\Theta = \R \times (0,\infty)$.
\end{enumerate}
\end{example}

\begin{aestheticbox}
Why bother with a model at all? Because data alone cannot answer a question
about the world; only data \emph{plus} an assumption about how the world
generated it can. The model is the bridge from ``what I saw'' to ``what is
true.'' Choose the bridge badly and even infinite data leads you to a confident
wrong answer. Every estimate in this chapter is conditional on its model being
roughly right -- a humility worth carrying.
\end{aestheticbox}

\section{The Method of Moments}
\label{ch30:mom}

The oldest estimation idea is also the most intuitive, and you have been using
its baby version your whole life. If you want to know the average height of
adults, you measure some adults and average their heights. You are matching a
\emph{sample average} to a \emph{population average}. The method of moments,
introduced by Karl Pearson in the $1890$s, simply industrializes this instinct.

\begin{intuitionbox}
The model tells you what the theoretical moments -- the mean $\E[X]$, the mean
square $\E[X^2]$, and so on -- ought to be, \emph{as functions of the unknown
parameters}. The data hand you empirical versions of those same moments. The
method of moments sets theory equal to data and solves for the parameters. It
is the statistical equivalent of ``$n$ unknowns, $n$ equations.''
\end{intuitionbox}

Recall the $k$-th \emph{population moment} $\mu_k = \E[X^k]$ and define the
$k$-th \emph{sample moment} as the obvious empirical average
\[
  \hat\mu_k = \frac{1}{n} \sum_{i=1}^n X_i^k.
\]
The law of large numbers (a previous chapter) guarantees
$\hat\mu_k \pto \mu_k$: sample moments converge to population moments. That is
the whole engine. If $\mu_k$ depends on $\theta$, and the sample moment is close
to $\mu_k$, then solving the equation should land you near the true $\theta$.

\begin{definition}[Method-of-moments estimator]
\label{ch30:def-mom}
For a model with $k$ unknown parameters, write the first $k$ population moments
$\mu_1(\theta),\dots,\mu_k(\theta)$ as functions of $\theta$. The
\emph{method-of-moments estimator} $\hat\theta_{\mathrm{MM}}$ is the solution of
the $k$ equations
\[
  \mu_1(\theta) = \hat\mu_1,\quad \mu_2(\theta) = \hat\mu_2,\quad\dots,\quad
  \mu_k(\theta) = \hat\mu_k.
\]
\end{definition}

\begin{example}[MoM for the Poisson rate]
\label{ch30:ex-mom-pois}
For $\Pois(\lambda)$ we know $\mu_1 = \E[X] = \lambda$. One parameter, so one
equation: $\lambda = \hat\mu_1 = \bar X$. The method-of-moments estimate of the
call rate is just the sample mean,
$\hat\lambda_{\mathrm{MM}} = \bar X = \frac1n\sum_i X_i$. If $200$ minutes
produced $560$ calls, $\hat\lambda = 2.8$ calls per minute.
\end{example}

\begin{example}[MoM for the Normal]
\label{ch30:ex-mom-normal}
Two parameters $(\mu,\sigma^2)$, so two equations. The first two population
moments are $\mu_1 = \mu$ and $\mu_2 = \E[X^2] = \sigma^2 + \mu^2$. Setting them
equal to the sample moments:
\[
  \mu = \hat\mu_1 = \bar X,
  \qquad
  \sigma^2 + \mu^2 = \hat\mu_2 = \frac1n\sum_i X_i^2.
\]
Substituting the first into the second,
\[
  \hat\sigma^2_{\mathrm{MM}} = \frac1n\sum_i X_i^2 - \bar X^2
  = \frac1n\sum_i (X_i - \bar X)^2.
\]
So MoM hands you the sample mean and the sample variance (with the $1/n$
divisor) -- comfortingly familiar faces.
\end{example}

\begin{example}[MoM for the Uniform -- where intuition gets interesting]
\label{ch30:ex-mom-unif}
Let $X_i \sim \Unif(0,\theta)$, so the data are spread between $0$ and an
unknown ceiling $\theta$. The mean of a $\Unif(0,\theta)$ is $\theta/2$, so the
moment equation is $\theta/2 = \bar X$, giving
\[
  \hat\theta_{\mathrm{MM}} = 2\bar X.
\]
Here is a delicious trap. Suppose your sample is
$0.1,\ 0.4,\ 0.9$. Then $\bar X = 0.467$ and
$\hat\theta_{\mathrm{MM}} = 0.933$. But you \emph{observed} a value of $0.9$!
The estimated ceiling is barely above a point you actually saw, and a slightly
unluckier sample could give $\hat\theta < \max_i X_i$ -- an estimated ceiling
\emph{below} an observation, which is logically impossible under the model. The
method of moments does not know that $\theta$ must be at least as large as the
biggest observation. It only knows about averages. This is our first hint that
intuition-by-averaging, though seductive, can be improved upon.
Figure~\ref{ch30:fig-mom-unif} draws the situation.
\end{example}

\begin{figure}[H]
\centering
\begin{tikzpicture}
\begin{axis}[
  width=13cm, height=4.2cm,
  xlabel={observed value}, ylabel={},
  xmin=0, xmax=1.05, ymin=0, ymax=1.2,
  ytick=\empty, ymajorgrids=false,
  xtick={0,0.2,0.4,0.6,0.8,1.0},
  axis y line=none, axis x line=bottom,
  title={Method of moments can place the estimated ceiling \emph{below} an observation},
]
% data points
\addplot[only marks, mark=*, mark size=2.6pt, distblue]
  coordinates {(0.1,0.4)(0.4,0.4)(0.9,0.4)};
% labels
\node[anchor=south, distblue, font=\scriptsize] at (axis cs:0.1,0.45) {$0.1$};
\node[anchor=south, distblue, font=\scriptsize] at (axis cs:0.4,0.45) {$0.4$};
\node[anchor=south, distblue, font=\scriptsize] at (axis cs:0.9,0.45) {$0.9$};
% MoM estimate
\draw[targetgreen, very thick] (axis cs:0.9333,0) -- (axis cs:0.9333,0.95);
\node[anchor=south, targetgreen!60!black, font=\scriptsize, align=center]
  at (axis cs:0.9333,0.96) {$\hat\theta_{\mathrm{MM}}=2\bar X$\\$=0.933$};
% max observation
\draw[trapred, very thick, dashed] (axis cs:0.9,0) -- (axis cs:0.9,0.78);
\node[anchor=south east, trapred, font=\scriptsize] at (axis cs:0.9,0.62) {$\max_i x_i=0.9$};
\end{axis}
\end{tikzpicture}
\caption{Three observations (dots) from $\Unif(0,\theta)$. The method-of-moments
ceiling $2\bar X = 0.933$ sits barely above the largest observation $0.9$. A
slightly unluckier draw would push $2\bar X$ \emph{below} $0.9$ -- an estimated
ceiling lower than a value we actually saw, which the model forbids. Averaging
alone is blind to this constraint.}
\label{ch30:fig-mom-unif}
\end{figure}

\begin{trapbox}
The method of moments is fast, requires no calculus, and gives a sensible
starting guess for almost any model -- which is exactly why it is still used to
initialize harder optimizations. But it can produce nonsensical answers (a
ceiling below an observation, a negative variance, a probability above one)
because it only listens to a few summary numbers and ignores the finer
structure of the likelihood. Treat it as a quick sketch, not a final portrait.
\end{trapbox}

\begin{historybox}
Karl Pearson built the method of moments around $1894$ while fitting skewed
biological data -- crab measurements from the Bay of Naples -- that refused to
look Gaussian. He modeled them as mixtures and matched up to five moments,
grinding through the algebra by hand. For thirty years it was \emph{the} way to
fit distributions. Then a young R.\,A.\ Fisher argued, in a series of papers
from $1912$ to $1922$, that matching moments was wasteful: it threw away
information the data contained. His alternative, maximum likelihood, eventually
won the field. The friendly rivalry between Pearson's and Fisher's schools is
one of the founding dramas of modern statistics -- and personal: the two men
famously could not stand each other.
\end{historybox}

\section{The Likelihood Function}
\label{ch30:likelihood}

Fisher's complaint was that the method of moments listens to only a handful of
summaries. The likelihood listens to everything. The whole idea fits in two
words: \emph{run probability backwards}. Ordinarily
$\P(\text{data} \given \theta)$ treats $\theta$ as fixed
and asks how probable the data are. Fisher said: I have the data; let me hold
\emph{them} fixed and view that same number as a function of $\theta$. Now it
ranks the candidate parameter values by how well each would have predicted what
actually happened.

\begin{definition}[Likelihood and log-likelihood]
\label{ch30:def-likelihood}
Given observed data $x_1,\dots,x_n$ and a model with parameter $\theta$, the
\emph{likelihood function} is the joint density (or mass) of the data, read as a
function of $\theta$ with the data held fixed:
\[
  L(\theta) = L(\theta; x_1,\dots,x_n) = \prod_{i=1}^n f(x_i \given \theta)
\]
for i.i.d.\ data with density/mass $f$. Its logarithm, the \emph{log-likelihood},
is
\[
  \ell(\theta) = \log L(\theta) = \sum_{i=1}^n \log f(x_i \given \theta).
\]
\end{definition}

\begin{trapbox}
$L(\theta)$ is \emph{not} a probability distribution over $\theta$. It does not
integrate to one over $\Theta$, and ``$L(\theta)$'' is not ``the probability
that the parameter is $\theta$.'' The parameter is a fixed number, not a random
variable -- it has no distribution to speak of. (When you \emph{do} put a
distribution on $\theta$ you have crossed into Bayesian inference, the subject
of the next chapter, and the object you get is the posterior, not the
likelihood.) The likelihood is a \emph{ranking}, not a probability: it tells you
which $\theta$ values make the data look plausible, in relative terms only.
\end{trapbox}

\begin{whybox}
Why multiply the densities? Because the observations are independent, and
independence makes joint probabilities multiply. The probability of seeing
\emph{this} value and \emph{that} value and \emph{the other} value is the
product of the individual probabilities. Why take the log? Three reasons, all
practical. First, products of $n$ small numbers underflow to zero on a computer;
sums of logs do not. Second, sums are far easier to differentiate than products
-- and we are about to differentiate. Third, the log turns the i.i.d.\ product
into a \emph{sum} over data points, which is exactly the form the law of large
numbers and the central limit theorem love. Because $\log$ is strictly
increasing, the $\theta$ that maximizes $L$ also maximizes $\ell$, so we lose
nothing.
\end{whybox}

\section{Maximum Likelihood Estimation}
\label{ch30:mle}

\begin{definition}[Maximum likelihood estimator]
\label{ch30:def-mle}
The \emph{maximum likelihood estimator} (MLE) is the parameter value that makes
the observed data most probable:
\[
  \hat\theta_{\mathrm{MLE}} = \arg\max_{\theta \in \Theta} L(\theta)
                            = \arg\max_{\theta \in \Theta} \ell(\theta).
\]
\end{definition}

The recipe, when $\ell$ is smooth, is the familiar calculus drill: differentiate
$\ell$, set the derivative (the \emph{score}) to zero, solve, and check it is a
maximum.
\[
  \text{score: } \quad s(\theta) = \frac{\partial \ell}{\partial \theta} = 0.
\]
(In Section~\ref{ch30:fisher} we will use the word \emph{score} for the
derivative of a \emph{single} observation's log-density,
$\partial_\theta\log f(X\given\theta)$. The two usages are consistent: because
$\ell$ is a sum over observations, the full-data score is the sum of the
per-observation scores. We name both ``score'' as the field does, but the
distinction matters when we take variances.)

\begin{example}[The MLE of a coin's bias -- the canonical case]
\label{ch30:ex-mle-bern}
You toss a coin $n$ times and see $k$ heads. With $X_i\sim\Bern(p)$ the mass
function is $f(x\given p)=p^x(1-p)^{1-x}$, so the likelihood is
\[
  L(p) = \prod_{i=1}^n p^{x_i}(1-p)^{1-x_i} = p^{k}(1-p)^{n-k},
  \qquad k = \sum_i x_i.
\]
Take logs:
\[
  \ell(p) = k\log p + (n-k)\log(1-p).
\]
Differentiate and set to zero:
\[
  \ell'(p) = \frac{k}{p} - \frac{n-k}{1-p} = 0
  \;\Longrightarrow\;
  k(1-p) = (n-k)p
  \;\Longrightarrow\;
  k = np
  \;\Longrightarrow\;
  \boxed{\;\hat p_{\mathrm{MLE}} = \frac{k}{n}.\;}
\]
The maximum likelihood estimate of the coin's bias is the observed fraction of
heads. For $61$ heads in $100$, $\hat p = 0.61$. The MLE recovers the most
natural estimate there is -- but now you know \emph{why} it is natural: it
is the value that makes your particular run of $61$ heads the most probable
outcome it could have been.
\end{example}

Figure~\ref{ch30:fig-loglik-bern} shows the log-likelihood for this coin. Notice
its single smooth peak at $0.61$, and notice how the curve gets \emph{sharper}
as the sample grows -- a foreshadowing of Fisher information in
Section~\ref{ch30:fisher}.

\begin{figure}[H]
\centering
\begin{tikzpicture}
\begin{axis}[
  width=12cm, height=7cm,
  xlabel={$p$ (candidate bias)}, ylabel={log-likelihood $\ell(p)$ (shifted)},
  domain=0.30:0.92, samples=160,
  legend pos=south west, legend cell align=left,
  grid=both, grid style={gray!18},
  xmin=0.30, xmax=0.92, ymin=-40, ymax=4,
  title={Log-likelihood for ``$61\%$ heads'' at three sample sizes},
]
\addplot[distblue, very thick] {12.2*ln(x) + 7.8*ln(1-x) - (12.2*ln(0.61)+7.8*ln(0.39))};
\addlegendentry{$n=20$}
\addplot[densityred, very thick, dashed] {61*ln(x) + 39*ln(1-x) - (61*ln(0.61)+39*ln(0.39))};
\addlegendentry{$n=100$}
\addplot[targetgreen, very thick, dotted] {305*ln(x) + 195*ln(1-x) - (305*ln(0.61)+195*ln(0.39))};
\addlegendentry{$n=500$}
\draw[gray, thick, dashed] (axis cs:0.61,-40) -- (axis cs:0.61,2);
\node[anchor=south] at (axis cs:0.61,1.2) {\footnotesize peak at $\hat p=0.61$};
\end{axis}
\end{tikzpicture}
\caption{The log-likelihood for a coin showing $61\%$ heads, shifted to peak at
zero. All three curves peak at $\hat p = 0.61$ regardless of $n$, but more data
makes the peak dramatically sharper -- the data ``insist'' on $0.61$ more
strongly. That sharpness is precisely what makes a large sample informative.}
\label{ch30:fig-loglik-bern}
\end{figure}

\begin{example}[The MLE of the Poisson rate]
\label{ch30:ex-mle-pois}
With $X_i\sim\Pois(\lambda)$, $f(x\given\lambda)=e^{-\lambda}\lambda^x/x!$, so
\[
  \ell(\lambda) = \sum_i \left(-\lambda + x_i\log\lambda - \log x_i!\right)
  = -n\lambda + \Big(\sum_i x_i\Big)\log\lambda - \sum_i \log x_i!.
\]
The last term does not involve $\lambda$, so it falls away under
differentiation:
\[
  \ell'(\lambda) = -n + \frac{\sum_i x_i}{\lambda} = 0
  \;\Longrightarrow\;
  \hat\lambda_{\mathrm{MLE}} = \frac{\sum_i x_i}{n} = \bar X.
\]
Here the MLE and the method-of-moments estimate coincide -- both equal $\bar X$.
They will not always agree, but when the model is ``nice'' (an exponential
family, matched at its natural moments) they often do.
\end{example}

\begin{example}[The MLE of the Normal -- two parameters at once]
\label{ch30:ex-mle-normal}
For $X_i\sim\Normal(\mu,\sigma^2)$,
\[
  \ell(\mu,\sigma^2)
  = -\frac{n}{2}\log(2\pi) - \frac{n}{2}\log\sigma^2
    - \frac{1}{2\sigma^2}\sum_i (x_i-\mu)^2.
\]
Set the two partial derivatives to zero. With respect to $\mu$:
\[
  \frac{\partial\ell}{\partial\mu}
  = \frac{1}{\sigma^2}\sum_i (x_i-\mu) = 0
  \;\Longrightarrow\;
  \hat\mu = \bar X.
\]
With respect to $\sigma^2$ (treat $v=\sigma^2$ as the variable):
\[
  \frac{\partial\ell}{\partial v}
  = -\frac{n}{2v} + \frac{1}{2v^2}\sum_i (x_i-\mu)^2 = 0
  \;\Longrightarrow\;
  \hat\sigma^2 = \frac1n\sum_i (x_i-\bar X)^2.
\]
Again the MLE matches the method of moments. Hold onto that $1/n$ divisor: it
will turn out to be biased, and the fix is one of the most famous footnotes in
statistics (Section~\ref{ch30:bias}).
\end{example}

\begin{example}[The MLE of the Uniform ceiling -- calculus fails, and that is the lesson]
\label{ch30:ex-mle-unif}
Return to $X_i\sim\Unif(0,\theta)$. The density is $f(x\given\theta)=1/\theta$
for $0\le x\le\theta$ and $0$ otherwise. The likelihood is
\[
  L(\theta) = \prod_{i=1}^n \frac{1}{\theta}\,\ind\set{0\le x_i\le\theta}
            = \frac{1}{\theta^n}\,\ind\set{\theta \ge \max_i x_i}.
\]
Here differentiating is a trap. As a function of $\theta$, $1/\theta^n$ is
\emph{decreasing}, so we want $\theta$ as small as possible -- but the indicator
forbids $\theta$ from dropping below the largest observation. The likelihood is
therefore maximized at the boundary:
\[
  \boxed{\;\hat\theta_{\mathrm{MLE}} = \max_i X_i.\;}
\]
Compare this to the method-of-moments answer $2\bar X$. The MLE can never be
less than an observed value -- it respects the logic of the model. This is the
clearest single demonstration in the book that maximum likelihood ``uses more
of the data'': it knows the maximum matters, while moment-matching only ever saw
the average.
\end{example}

\begin{figure}[H]
\centering
\begin{tikzpicture}
\begin{axis}[
  width=12cm, height=6cm,
  xlabel={$\theta$ (candidate ceiling)}, ylabel={likelihood $L(\theta)$},
  domain=0.9:2.2, samples=120,
  xmin=0.3, xmax=2.2, ymin=0, ymax=2,
  grid=both, grid style={gray!18},
  title={Uniform likelihood with $\max_i x_i = 0.9$: maximized at the boundary},
]
\addplot[densityred, very thick] {1/x^4};
\addplot[draw=none, fill=gray!12] coordinates {(0.3,0)(0.9,0)(0.9,2)(0.3,2)};
\draw[distblue, very thick, dashed] (axis cs:0.9,0) -- (axis cs:0.9,1.524);
\node[anchor=west] at (axis cs:0.92,1.5) {\footnotesize $\hat\theta = \max_i x_i = 0.9$};
\node[anchor=west, gray!60!black] at (axis cs:0.42,1.0) {\footnotesize $L=0$ here};
\node[anchor=west, gray!60!black] at (axis cs:0.42,0.7) {\footnotesize ($\theta<\max x_i$)};
\end{axis}
\end{tikzpicture}
\caption{For $\Unif(0,\theta)$ with $n=4$ and largest observation $0.9$, the
likelihood is zero whenever $\theta<0.9$ (an impossible ceiling) and decays as
$\theta^{-4}$ above it. The maximum sits exactly at the boundary
$\hat\theta=0.9$. No derivative is set to zero -- the optimum is a corner.}
\label{ch30:fig-unif-lik}
\end{figure}

\begin{happeningbox}
What is maximum likelihood \emph{really} doing? Picture the parameter space as a
landscape and the log-likelihood as its elevation. Each data point tilts the
landscape a little, pushing the summit toward parameter values that explain that
point well. With i.i.d.\ data these little tilts \emph{add up} (because logs
turn the product into a sum), so the summit settles at the consensus of all the
data points. As $n$ grows, more votes accumulate and the summit grows sharper
and more decisively located -- the law of large numbers acting on the score
function. Deeper still: maximizing the likelihood is the same as minimizing the
Kullback--Leibler divergence from the true distribution to your model. The MLE
is the member of your model family that is, in an information-theoretic sense,
\emph{closest to the truth} given the evidence. We make this precise in the
information-theory chapter; here, just feel the shape of it.
\end{happeningbox}

\begin{connectionbox}
Maximum likelihood is an optimization problem, so the whole machinery of
calculus comes to bear: the first-order condition (score $=0$) is just
``gradient equals zero,'' and checking it is a maximum uses the second
derivative being negative. For multi-parameter models the score becomes a
gradient \emph{vector} and the curvature becomes a Hessian \emph{matrix} -- so
linear algebra enters too, and the matrix of expected curvatures is the Fisher
information matrix of Section~\ref{ch30:fisher}. When the score equation has no
closed form (logistic regression, mixtures, most of modern machine learning),
you climb the likelihood numerically with gradient ascent or Newton's method --
literally the same algorithms that train neural networks. See the companion
calculus and linear algebra readers for the optimization background.
\end{connectionbox}

\subsection{The invariance property}
\label{ch30:invariance}

Maximum likelihood has a property the method of moments lacks, and it saves real
work.

\begin{theorem}[Invariance of the MLE]
\label{ch30:thm-invariance}
If $\hat\theta$ is the MLE of $\theta$ and $g$ is any function, then the MLE of
$g(\theta)$ is $g(\hat\theta)$.
\end{theorem}

\begin{whybox}
The reason is almost a tautology, which is why it is so reliable. The MLE is the
$\theta$ that sits at the top of the likelihood hill. If you re-label the
horizontal axis by a transformation $g$, the hill does not move -- the same
\emph{point} is still the summit, it just wears a new name $g(\hat\theta)$.
Maximization cares about \emph{where} the peak is, not what you call the
location.
\end{whybox}

\begin{example}[Invariance in action]
\label{ch30:ex-invariance}
The MLE of the Bernoulli parameter is $\hat p = k/n$. Suppose you actually care
about the \emph{odds} $g(p)=p/(1-p)$. By invariance, no new calculus is needed:
the MLE of the odds is simply $\hat p/(1-\hat p) = k/(n-k)$. Likewise, the MLE
of the Normal standard deviation $\sigma$ is $\sqrt{\hat\sigma^2}$, the square
root of the MLE of the variance. You never have to re-derive anything.
\end{example}

\section{Judging an Estimator: Bias, Variance, Consistency}
\label{ch30:bias}

We now have two machines for producing estimates. But producing an estimate is
not the same as producing a \emph{good} estimate. How do we grade an estimator?
The estimator is random, so we grade its sampling distribution. Three properties
matter most.

\begin{definition}[Bias]
\label{ch30:def-bias}
The \emph{bias} of an estimator $\hat\theta$ of $\theta$ is
\[
  \operatorname{Bias}(\hat\theta) = \E[\hat\theta] - \theta.
\]
The estimator is \emph{unbiased} if its bias is zero, i.e.\
$\E[\hat\theta]=\theta$ for every $\theta$: on average, across many hypothetical
repetitions of the experiment, it lands on the truth.
\end{definition}

\begin{intuitionbox}
Bias is a systematic error -- a scope that is sighted off-center. Even with
infinite patience (infinitely many repeated experiments), the average of a
biased estimator misses the bullseye. Unbiasedness says: \emph{no} systematic
lean. It does not say any single estimate is correct, only that the errors
cancel on average.
\end{intuitionbox}

\begin{example}[The sample mean is unbiased]
\label{ch30:ex-unbiased-mean}
If $X_i$ have mean $\mu$ then by linearity of expectation,
\[
  \E[\bar X] = \E\!\left[\frac1n\sum_i X_i\right]
             = \frac1n\sum_i \E[X_i] = \frac1n\cdot n\mu = \mu.
\]
The sample mean is unbiased for the population mean, full stop -- it needs no
assumption beyond a finite mean. This is why $\bar X$ is the workhorse of
statistics.
\end{example}

\begin{example}[The MLE of the variance is biased -- and the famous fix]
\label{ch30:ex-biased-var}
Recall the MLE $\hat\sigma^2 = \frac1n\sum_i(X_i-\bar X)^2$. Is it unbiased? Let
us compute its expectation. Write $\sigma^2$ for the true variance. A standard
identity expands the sum of squared deviations:
\[
  \sum_i (X_i-\bar X)^2 = \sum_i (X_i-\mu)^2 - n(\bar X-\mu)^2.
\]
Take expectations. Each $\E[(X_i-\mu)^2]=\sigma^2$, giving $n\sigma^2$ for the
first term. And $\E[(\bar X-\mu)^2]=\Var(\bar X)=\sigma^2/n$, so the second term
contributes $n\cdot\sigma^2/n=\sigma^2$. Hence
\[
  \E\!\left[\sum_i(X_i-\bar X)^2\right] = n\sigma^2 - \sigma^2 = (n-1)\sigma^2,
\]
and therefore
\[
  \E[\hat\sigma^2] = \frac{n-1}{n}\,\sigma^2 < \sigma^2.
\]
The MLE \emph{systematically underestimates} the variance. The cure is to divide
by $n-1$ instead of $n$:
\[
  S^2 = \frac{1}{n-1}\sum_i (X_i-\bar X)^2,
  \qquad \E[S^2] = \sigma^2.
\]
This $S^2$ is the \emph{sample variance} you have seen everywhere with its
mysterious $n-1$. Now the mystery is gone.
\end{example}

\begin{whybox}
Why $n-1$ and not $n$? Because you used the data twice. To measure spread around
the center you first had to \emph{estimate} the center with $\bar X$ -- and
$\bar X$ is, by construction, the point that sits closest to your particular
sample. Deviations measured from $\bar X$ are therefore a touch too small,
because $\bar X$ was pulled toward the data. You have spent one ``degree of
freedom'' locating the center, leaving only $n-1$ independent pieces of
information about spread. Dividing by $n-1$ compensates exactly. The phrase
\emph{degrees of freedom} is just bookkeeping for ``how many independent
residuals are really left.''
\end{whybox}

\begin{trapbox}
Maximum likelihood does \emph{not} promise unbiasedness, and the variance MLE is
the standard counterexample. The MLE's virtues are large-sample virtues
(consistency, efficiency); in small samples it can lean. Do not assume ``MLE''
means ``unbiased.'' Conversely, do not over-worship unbiasedness either -- as we
are about to see, a little bias can be a bargain.
\end{trapbox}

\begin{definition}[Variance and mean squared error]
\label{ch30:def-mse}
The \emph{variance} of an estimator is $\Var(\hat\theta)$, the spread of its
sampling distribution. The \emph{mean squared error} combines both kinds of
error:
\[
  \mathrm{MSE}(\hat\theta) = \E\big[(\hat\theta-\theta)^2\big]
  = \underbrace{\Var(\hat\theta)}_{\text{wobble}}
  + \underbrace{\operatorname{Bias}(\hat\theta)^2}_{\text{lean}}.
\]
\end{definition}

\begin{whybox}
The MSE decomposition is just the ``add and subtract the mean'' trick. Write
$\hat\theta-\theta = (\hat\theta-\E\hat\theta)+(\E\hat\theta-\theta)$. Square it:
the first piece squared has expectation $\Var(\hat\theta)$, the second piece is
the constant $\operatorname{Bias}^2$, and the cross term is
$2\,\E[\hat\theta-\E\hat\theta]\cdot\operatorname{Bias}=0$ because the first
factor has mean zero. So total expected squared error cleanly splits into
\emph{wobble plus lean}. This little identity is the seed of the entire
bias--variance tradeoff.
\end{whybox}

\begin{definition}[Consistency]
\label{ch30:def-consistency}
An estimator is \emph{consistent} if it converges in probability to the truth as
data accumulate:
$\hat\theta_n \pto \theta$, meaning for every $\eps>0$,
$\P(\abs{\hat\theta_n - \theta} > \eps)\to 0$ as $n\to\infty$.
\end{definition}

\begin{intuitionbox}
Consistency is the bare-minimum sanity check: \emph{if you collect enough data,
you get the right answer.} An inconsistent estimator is broken at the
foundation -- no amount of data rescues it. A clean sufficient condition: if
$\operatorname{Bias}(\hat\theta_n)\to 0$ \emph{and}
$\Var(\hat\theta_n)\to 0$, then $\mathrm{MSE}\to 0$, and convergence in mean
square forces convergence in probability (Chebyshev). So \emph{shrinking bias
plus shrinking variance} $\Rightarrow$ consistency.
\end{intuitionbox}

\begin{example}[Both variance estimators are consistent]
\label{ch30:ex-consistent-var}
The biased MLE $\hat\sigma^2$ has bias $-\sigma^2/n \to 0$ and variance of order
$1/n \to 0$, so it is consistent despite being biased. The unbiased $S^2$ is
also consistent. As $n\to\infty$ the factor $(n-1)/n\to 1$ and the two estimators
become indistinguishable. \emph{The $n-1$ versus $n$ debate is a small-sample
debate}; asymptotically it evaporates. This is the recurring moral of the
chapter: bias usually fades with data, so the real long-run battle is over
variance.
\end{example}

\begin{figure}[H]
\centering
\begin{tikzpicture}
\begin{axis}[
  width=13cm, height=6.2cm,
  xlabel={target}, ylabel={frequency},
  ymin=0, ymax=1.05, xmin=-3.2, xmax=3.2,
  ytick=\empty,
  grid=both, grid style={gray!15},
  title={Four estimators as sampling distributions (truth at $0$)},
  domain=-3.2:3.2, samples=160,
  legend pos=north west, legend style={font=\footnotesize}, legend cell align=left,
]
\addplot[targetgreen, very thick] {exp(-x^2/(2*0.16))};
\addlegendentry{unbiased, low variance (best)}
\addplot[distblue, very thick, dashed] {0.40*exp(-x^2/(2*1.0))};
\addlegendentry{unbiased, high variance}
\addplot[quotientorange, very thick, dotted] {exp(-(x-1.1)^2/(2*0.16))};
\addlegendentry{biased, low variance}
\addplot[densityred, very thick, dashdotted] {0.40*exp(-(x-1.1)^2/(2*1.0))};
\addlegendentry{biased, high variance (worst)}
\draw[gray, thick] (axis cs:0,0) -- (axis cs:0,1.05);
\node[anchor=south, gray!50!black] at (axis cs:0,1.0) {\footnotesize truth};
\end{axis}
\end{tikzpicture}
\caption{Bias is where the distribution is \emph{centered} relative to the truth
(horizontal shift); variance is how \emph{wide} it is. A good estimator is both
centered on the truth and tight around it. The mean squared error is the average
squared distance from the truth, capturing both the shift and the width at once.}
\label{ch30:fig-bias-variance}
\end{figure}

\section{The Bias--Variance Tradeoff}
\label{ch30:bv-tradeoff}

Decompose MSE into wobble plus lean and a strategic possibility leaps out: maybe
you can \emph{accept} a little lean (bias) in exchange for far less wobble
(variance), and come out ahead on total error. Unbiasedness is not sacred;
\emph{small MSE} is the real goal.

\begin{example}[Deliberately biasing the coin estimate]
\label{ch30:ex-shrinkage}
The unbiased estimate of a coin's bias is $\hat p = k/n$. But near the extremes
it behaves badly: with $0$ heads in $5$ tosses it shouts ``$p=0$, this coin
\emph{never} lands heads,'' a wildly overconfident claim from five tosses.
Consider instead the gently shrunk estimator
\[
  \tilde p = \frac{k+1}{n+2},
\]
which nudges every estimate toward $1/2$ (it is, not coincidentally, the
Bayesian posterior mean under a uniform prior -- the next chapter). With $0$
heads in $5$ it returns $1/7\approx 0.14$ instead of $0$: still ``probably
tail-heavy,'' but not absurdly certain. $\tilde p$ is \emph{biased} -- it leans
toward $1/2$ -- yet for moderate $n$ and $p$ not too near $0$ or $1$ it often has
\emph{smaller MSE} than the unbiased $\hat p$, because the variance reduction
near the boundaries outweighs the small bias. A pinch of bias bought a calmer,
more reliable estimator.
\end{example}

\begin{applicationbox}[The bias--variance tradeoff is the soul of machine learning]
This single idea -- trade a little bias for a lot less variance -- organizes the
entire practice of supervised learning. A model's expected prediction error
decomposes (exactly as our MSE did) into
\[
  \text{error} = \text{bias}^2 + \text{variance} + \text{irreducible noise}.
\]
\begin{itemize}
\item A \emph{too-simple} model (a straight line for a curvy world) has high
  bias: it cannot represent the truth no matter how much data you give it. It
  \emph{underfits}.
\item A \emph{too-flexible} model (a degree-$30$ polynomial, a giant decision
  tree) has high variance: it contorts itself to chase the noise in this
  particular training set and changes wildly with a new sample. It
  \emph{overfits}.
\end{itemize}
Every regularization knob in machine learning -- ridge and lasso penalties,
shrinkage, dropout, early stopping, pruning, limiting tree depth -- is
\emph{deliberately introduced bias} that buys a large reduction in variance.
Ridge regression, for instance, shrinks coefficients toward zero exactly as
$\tilde p$ shrank toward $1/2$: biased, but lower MSE on fresh data. When a
practitioner says a model is overfitting, they are saying its variance is too
high; the cure is almost always to add bias. Much of applied machine learning is
the humble formula $\mathrm{MSE}=\Var+\operatorname{Bias}^2$ worked out in
practice.
\end{applicationbox}

\begin{paradoxbox}[Stein's paradox: the sample mean is not always best]
Here is one of the most startling results in all of statistics, and it lives
exactly at the bias--variance frontier. Suppose you estimate three or more
\emph{unrelated} means at once -- say the true batting averages of several
players, or the means of three independent $\Normal(\mu_i,1)$ variables. The
obvious estimator is to use each sample mean for its own parameter, and each
sample mean is individually unbiased and efficient. Yet Charles Stein proved in
$1956$ that for dimension $\ge 3$ this obvious estimator is \emph{inadmissible}:
the James--Stein estimator, which deliberately \emph{shrinks} all the sample
means toward a common point, has \emph{strictly smaller total MSE} for every
true value of the parameters. Borrowing strength across problems that have
nothing to do with each other beats handling each alone. The resolution is the
chapter's thesis taken to its extreme: a dose of bias (shrinkage) can buy enough
variance reduction to win even when the unbiased estimator was already efficient.
Modern hierarchical and regularized models -- and much of machine learning's
success -- are Stein's insight industrialized.
\end{paradoxbox}

\begin{figure}[H]
\centering
\begin{tikzpicture}
\begin{axis}[
  width=13cm, height=6.2cm,
  xlabel={model complexity (flexibility)}, ylabel={expected test error},
  domain=0.4:5, samples=120,
  xmin=0.4, xmax=5, ymin=0, ymax=6,
  grid=both, grid style={gray!15},
  legend pos=north east, legend style={font=\footnotesize}, legend cell align=left,
  xtick=\empty,
  title={The bias--variance tradeoff: total error is U-shaped},
]
\addplot[distblue, very thick] {9/(x^2)+0.3};
\addlegendentry{$\text{bias}^2$ (falls with complexity)}
\addplot[densityred, very thick, dashed] {0.18*x^2};
\addlegendentry{variance (rises with complexity)}
\addplot[targetgreen, ultra thick] {9/(x^2)+0.3+0.18*x^2};
\addlegendentry{total error (U-shaped)}
\draw[gray, thick, dotted] (axis cs:2.46,0) -- (axis cs:2.46,2.6);
\node[anchor=south, gray!55!black] at (axis cs:2.46,2.6) {\footnotesize sweet spot};
\node[anchor=west, distblue!70!black] at (axis cs:0.55,4.6) {\footnotesize underfit};
\node[anchor=east, densityred!70!black] at (axis cs:4.9,5.2) {\footnotesize overfit};
\end{axis}
\end{tikzpicture}
\caption{As model flexibility rises, squared bias falls but variance climbs.
Their sum is U-shaped, and the best model lives at the bottom of the U -- not the
simplest, not the most complex, but the one that balances lean against wobble.
Choosing that point (via cross-validation) is most of applied modeling.}
\label{ch30:fig-bv-curve}
\end{figure}

\section{Fisher Information and the Cramer--Rao Bound}
\label{ch30:fisher}

We have seen that variance is the long-run battleground: bias usually fades, but
how small can variance get? Remarkably, there is a \emph{floor} -- a hard limit,
set by the model itself, below which no unbiased estimator's variance can drop.
The floor is governed by a single quantity that measures how much a single
observation can possibly tell you about the parameter: the \emph{Fisher
information}.

\begin{intuitionbox}
Look again at Figure~\ref{ch30:fig-loglik-bern}. With lots of data the
log-likelihood is a \emph{sharp} peak; with little data it is a \emph{gentle}
mound. A sharp peak means the data discriminate strongly between nearby
parameter values -- shift $\theta$ a hair and the likelihood plummets, so you
know $\theta$ precisely. A gentle mound means many $\theta$ values explain the
data nearly equally well, so you are uncertain. \emph{Information is curvature.}
Fisher information measures how sharply curved the log-likelihood is at its peak,
on average.
\end{intuitionbox}

\begin{definition}[Fisher information]
\label{ch30:def-fisher}
For a single observation $X$ with log-density $\log f(X\given\theta)$, the
\emph{score} is $s(\theta)=\partial_\theta \log f(X\given\theta)$ and the
\emph{Fisher information} is its variance,
\[
  I(\theta) = \Var\big(s(\theta)\big)
            = \E\!\left[\left(\frac{\partial}{\partial\theta}\log f(X\given\theta)\right)^{\!2}\right].
\]
Under mild regularity conditions this equals the expected curvature,
\[
  I(\theta) = -\,\E\!\left[\frac{\partial^2}{\partial\theta^2}\log f(X\given\theta)\right].
\]
For $n$ i.i.d.\ observations the information adds: $I_n(\theta)=n\,I(\theta)$.
\end{definition}

\begin{whybox}
Two facts make this definition click. First, the score has \emph{mean zero}:
$\E[s(\theta)]=\int (\partial_\theta f/f)\,f\,\dd x
= \partial_\theta\!\int f\,\dd x = \partial_\theta 1 = 0$. So the variance of the
score is just $\E[s^2]$ -- the first formula. Second, differentiating the
identity $\E[s]=0$ once more and applying the product rule yields
$\E[\partial_\theta s] + \E[s^2]=0$, i.e.\
$\E[s^2] = -\E[\partial_\theta s] = -\E[\partial_\theta^2 \log f]$ -- the second
formula. So ``variance of the score'' and ``expected curvature of the
log-likelihood'' are the same number, viewed from two sides. And information
\emph{adds} over independent observations because the total score is a sum of
independent mean-zero scores, and variances of independent things add: every
fresh data point contributes its own $I(\theta)$ to the pile. That additivity is
why ``more data'' so reliably means ``more precision.''
\end{whybox}

\begin{theorem}[Cramer--Rao lower bound]
\label{ch30:thm-crlb}
Under regularity conditions, \emph{any} unbiased estimator $\hat\theta$ of
$\theta$ based on $n$ i.i.d.\ observations satisfies
\[
  \Var(\hat\theta) \;\ge\; \frac{1}{n\,I(\theta)}.
\]
No unbiased estimator can do better than the inverse of the total Fisher
information. An estimator attaining the bound is called \emph{efficient}.
\end{theorem}

\begin{happeningbox}
The Cramer--Rao bound says that precision is a finite resource handed out by the
model, not by your cleverness. The model fixes $I(\theta)$ -- how much one
observation can reveal -- and then \emph{no} unbiased procedure, however
ingenious, can extract more than $nI(\theta)$ worth of certainty. Two
consequences worth absorbing. First, variance falls like $1/n$: to halve your
standard error you must \emph{quadruple} your sample. That is the iron $\sqrt n$
law behind every ``margin of error'' and the reason large studies are
expensive. Second, the bound is a \emph{benchmark}: compute $1/(nI)$, compare it
to your estimator's actual variance, and you know at once whether your estimator
is wasting information or squeezing the model dry.
\end{happeningbox}

\begin{example}[The CRLB for a coin, and why $\hat p=k/n$ is efficient]
\label{ch30:ex-crlb-bern}
For $\Bern(p)$, $\log f(x\given p)=x\log p+(1-x)\log(1-p)$. Differentiate twice:
\[
  \partial_p^2 \log f = -\frac{x}{p^2} - \frac{1-x}{(1-p)^2}.
\]
Take $-\E[\cdot]$ using $\E[X]=p$:
\[
  I(p) = \frac{p}{p^2} + \frac{1-p}{(1-p)^2}
       = \frac1p + \frac1{1-p}
       = \frac{1}{p(1-p)}.
\]
So the Cramer--Rao floor is
\[
  \Var(\hat p) \ge \frac{1}{nI(p)} = \frac{p(1-p)}{n}.
\]
Now check the actual MLE $\hat p=k/n=\bar X$. Since $k\sim\Bin(n,p)$,
$\Var(\hat p) = \Var(k)/n^2 = np(1-p)/n^2 = p(1-p)/n$. It \emph{exactly} attains
the bound. The natural ``fraction of heads'' is not merely reasonable -- it is
\emph{efficient}: no unbiased estimator of $p$ can ever beat it. Notice too that
$I(p)=1/[p(1-p)]$ blows up as $p\to 0$ or $1$: a near-deterministic coin is, per
toss, enormously informative when it does surprise you, which is why rare-event
estimation behaves so differently from coin-flip estimation.
\end{example}

\begin{example}[CRLB for the Poisson rate]
\label{ch30:ex-crlb-pois}
For $\Pois(\lambda)$, $\log f = -\lambda + x\log\lambda - \log x!$, so
$\partial_\lambda^2 \log f = -x/\lambda^2$ and
$I(\lambda) = -\E[-X/\lambda^2] = \lambda/\lambda^2 = 1/\lambda$. The floor is
$\Var(\hat\lambda)\ge \lambda/n$. The MLE $\hat\lambda=\bar X$ has variance
$\Var(\bar X)=\lambda/n$ (since $\Var(X)=\lambda$ for a Poisson). Efficient
again. Maximum likelihood keeps hitting the floor -- and that is no accident.
\end{example}

\begin{aestheticbox}[Why maximum likelihood is the default]
The examples are not coincidences. Under regularity conditions the MLE is, for
large $n$, \emph{asymptotically} the best estimator possible:
\[
  \hat\theta_{\mathrm{MLE}} \;\approx\; \Normal\!\left(\theta,\; \frac{1}{nI(\theta)}\right)
  \quad\text{for large } n.
\]
In words: it is asymptotically unbiased, asymptotically normal, and its variance
asymptotically attains the Cramer--Rao floor (asymptotic efficiency). That is
three things at once -- it lands on the truth, its errors are bell-shaped (so you
can build confidence intervals, two chapters from now), and it wrings every drop
of information from the model. This is why maximum likelihood is the default
estimation method across statistics, econometrics, and machine learning: in the
large-sample regime it is, provably, hard to beat. Fisher's thirty-year argument
with Pearson ended here.
\end{aestheticbox}

\begin{paradoxbox}[When maximum likelihood goes wrong]
Maximum likelihood is the hero of this chapter, so it is only fair to show its
scars; ``asymptotically optimal'' hides real failures.
\begin{itemize}
\item \textbf{The unbounded-likelihood catastrophe.} Fit a mixture of two
  Normals by maximum likelihood. Center one component exactly on a single data
  point and shrink its variance toward zero: that point's density blows up, and
  $L\to\infty$. The global maximum is a useless spike, not the sensible fit.
  The MLE, taken literally, does not exist.
\item \textbf{Too many parameters (Neyman--Scott).} Suppose each pair of
  observations has its own mean but all share a common variance $\sigma^2$. As
  the number of pairs grows, the MLE of $\sigma^2$ converges to $\sigma^2/2$ --
  consistently \emph{wrong} by a factor of two. When parameters proliferate as
  fast as data, the large-sample guarantees collapse.
\item \textbf{Boundary trouble.} For $\Unif(0,\theta)$ the MLE $\max_i X_i$ is
  biased (it is always below $\theta$) and not asymptotically normal -- its error
  shrinks like $1/n$, faster than the usual $1/\sqrt n$, and follows an
  exponential, not a bell. The clean asymptotic theory simply does not apply at
  a boundary.
\end{itemize}
The resolution is not to abandon maximum likelihood but to respect its fine
print: the optimality theorems require a smooth, bounded likelihood, a fixed
number of parameters, and an interior optimum. When those hold -- which is most
of the time -- the MLE is excellent. When they fail, reach for penalization,
Bayesian priors (next chapter), or a bias correction. A method this useful is
worth knowing at its edges, not just in its comfort zone.
\end{paradoxbox}

\begin{figure}[H]
\centering
\begin{tikzpicture}
\begin{axis}[
  width=12cm, height=6cm,
  xlabel={$p$}, ylabel={Fisher information $I(p)=1/[p(1-p)]$},
  domain=0.06:0.94, samples=160,
  xmin=0, xmax=1, ymin=0, ymax=20,
  grid=both, grid style={gray!18},
  title={Information per coin toss is lowest at $p=1/2$},
]
\addplot[distblue, very thick] {1/(x*(1-x))};
\draw[gray, dashed] (axis cs:0.5,0) -- (axis cs:0.5,4);
\node[anchor=south] at (axis cs:0.5,4) {\footnotesize min $I=4$ at $p=1/2$};
\end{axis}
\end{tikzpicture}
\caption{Fisher information for a single Bernoulli toss. It is \emph{smallest} at
$p=1/2$ (a fair coin is the hardest bias to pin down -- each toss is maximally
unpredictable, so individual outcomes say least) and rises steeply near the
extremes, where the rare surprising outcome is very telling. Each toss therefore
varies in how much it teaches you, depending on where the truth lies.}
\label{ch30:fig-fisher-bern}
\end{figure}

\section{Putting It Together: A Full Worked Estimation}
\label{ch30:worked}

Here is one example carried end to end, the way you would handle it in practice,
so you can watch all the pieces interlock.

\begin{example}[Estimating an exponential rate from waiting times]
\label{ch30:ex-full-expo}
A web server logs the gaps between requests (in seconds):
\[
  0.5,\quad 1.2,\quad 0.3,\quad 2.1,\quad 0.9,\quad 1.5.
\]
Model them as i.i.d.\ $\Expo(\lambda)$, density $f(x\given\lambda)=\lambda
e^{-\lambda x}$ for $x>0$, where $\lambda$ is the request \emph{rate}. The sum is
$\sum x_i = 6.5$, so $\bar x = 6.5/6 \approx 1.083$.

\textbf{Step 1: Method of moments.} $\E[X]=1/\lambda$, so the moment equation is
$1/\lambda=\bar x$, giving $\hat\lambda_{\mathrm{MM}}=1/\bar x \approx 0.923$
requests per second.

\textbf{Step 2: Maximum likelihood.} The log-likelihood is
\[
  \ell(\lambda) = \sum_i (\log\lambda - \lambda x_i)
                = n\log\lambda - \lambda\sum_i x_i.
\]
Differentiate and solve:
\[
  \ell'(\lambda) = \frac n\lambda - \sum_i x_i = 0
  \;\Longrightarrow\;
  \hat\lambda_{\mathrm{MLE}} = \frac{n}{\sum_i x_i} = \frac{1}{\bar x}
  \approx 0.923.
\]
For the exponential, MoM and MLE coincide -- both give the reciprocal of the
mean gap. (By invariance, the MLE of the \emph{mean} gap $1/\lambda$ is just
$\bar x\approx 1.083$ seconds.)

\textbf{Step 3: Fisher information and the precision floor.} With
$\log f = \log\lambda - \lambda x$ we get $\partial_\lambda^2\log f =
-1/\lambda^2$, so $I(\lambda)=1/\lambda^2$ and the Cramer--Rao floor is
\[
  \Var(\hat\lambda) \ge \frac{1}{nI(\lambda)} = \frac{\lambda^2}{n}.
\]
Plugging in $\hat\lambda\approx0.923$ and $n=6$ gives an approximate standard
error $\sqrt{\lambda^2/n}=\lambda/\sqrt n \approx 0.923/2.449 \approx 0.38$. So
our estimate is roughly $0.92 \pm 0.38$ requests/second -- usefully informative,
but with only six observations, not precise. To halve that $0.38$ we would need
about $24$ observations: the $\sqrt n$ law again.

\textbf{Step 4: Bias check.} It turns out $1/\bar x$ is slightly biased upward
for small $n$ (by Jensen's inequality, $\E[1/\bar X] > 1/\E[\bar X]=\lambda$); a
common correction is $\hat\lambda_{\text{corrected}}=\frac{n-1}{n}\cdot
\frac{1}{\bar x}$. With $n=6$ the correction multiplies by $5/6$, giving
$\approx 0.77$. As $n\to\infty$ the correction vanishes -- the same moral one
last time: \emph{bias fades, variance is forever}.
\end{example}

Figure~\ref{ch30:fig-sampling-narrow} previews what the lab below will let you
see for yourself: the sampling distribution of $\hat\lambda=\bar X$ is centred on
the truth and tightens as $n$ grows.

\begin{figure}[H]
\centering
\begin{tikzpicture}
\begin{axis}[
  width=13cm, height=6cm,
  xlabel={estimate $\hat\lambda=\bar X$}, ylabel={sampling density},
  domain=2.0:4.0, samples=160,
  xmin=2.0, xmax=4.0, ymin=0, ymax=3.6,
  ytick=\empty,
  grid=both, grid style={gray!15},
  legend pos=north west, legend style={font=\footnotesize}, legend cell align=left,
  title={Sampling distribution of $\hat\lambda$ for $\Pois(\lambda=3)$ narrows with $n$},
]
% N(3, 3/n): sd = sqrt(3/n). Normalized peaks.
\addplot[distblue, very thick]
  {1/(sqrt(2*pi*3/50))*exp(-(x-3)^2/(2*3/50))};
\addlegendentry{$n=50$ (SE $\approx 0.245$)}
\addplot[targetgreen, very thick, dashed]
  {1/(sqrt(2*pi*3/200))*exp(-(x-3)^2/(2*3/200))};
\addlegendentry{$n=200$ (SE $\approx 0.122$)}
\draw[gray, thick, dashed] (axis cs:3,0) -- (axis cs:3,3.5);
\node[anchor=south, gray!55!black] at (axis cs:3,3.5) {\footnotesize truth $\lambda=3$};
\end{axis}
\end{tikzpicture}
\caption{Two sampling distributions of the Poisson MLE $\hat\lambda=\bar X$, by
the asymptotic-normality approximation $\hat\lambda\approx\Normal(\lambda,
\lambda/n)$. Both centre on the truth $3$ (consistency); quadrupling $n$ from
$50$ to $200$ halves the standard error from $\approx 0.245$ to $\approx 0.122$
(the $\sqrt n$ law) and so doubles the peak height. The shaded bell shape is
exactly what the histogram in the lab below should reproduce.}
\label{ch30:fig-sampling-narrow}
\end{figure}

\begin{labbox}
Before the exercises, simulate the whole story yourself. Draw $1000$ datasets of
size $n=50$ from $\Pois(\lambda=3)$. For each, compute $\hat\lambda=\bar X$.
Histogram the $1000$ estimates: it should look bell-shaped (asymptotic
normality), centered near $3$ (consistency), with spread close to
$\sqrt{\lambda/n}=\sqrt{3/50}\approx 0.245$ (the Cramer--Rao floor). Then redo it
with $n=200$ and watch the histogram narrow by a factor of $2$ as $n$ quadruples.
You will have \emph{seen}, in one afternoon, every theorem in this chapter
breathe.
\end{labbox}

\section{Exercises}
\label{ch30:exercises}

\begin{exercise}[Geometric MLE]
\label{ch30:exc-geom}
Let $X_1,\dots,X_n \sim \Geom(p)$ on $\set{1,2,3,\dots}$ with
$\P(X=k)=(1-p)^{k-1}p$. Write down the log-likelihood, and show the MLE is
$\hat p = 1/\bar X$. Interpret: why is the estimated success probability the
reciprocal of the average number of trials?
\end{exercise}

\begin{exercise}[Method of moments vs.\ MLE for the uniform]
\label{ch30:exc-unif}
For $X_i\sim\Unif(0,\theta)$ you found $\hat\theta_{\mathrm{MM}}=2\bar X$ and
$\hat\theta_{\mathrm{MLE}}=\max_i X_i$.
(a) Show $\E[2\bar X]=\theta$, so the MoM estimator is unbiased.
(b) It is a fact that $\E[\max_i X_i]=\frac{n}{n+1}\theta$; find the bias of the
MLE and propose a multiplicative correction making it unbiased.
(c) Given the data $0.1, 0.4, 0.9$, compute both estimates and explain which one
you would trust and why.
\end{exercise}

\begin{exercise}[The $n-1$ in your bones]
\label{ch30:exc-n-minus-1}
Using $\E[(\bar X-\mu)^2]=\sigma^2/n$, re-derive
$\E\big[\frac1n\sum_i(X_i-\bar X)^2\big]=\frac{n-1}{n}\sigma^2$ from scratch, and
explain in one sentence, without formulas, why dividing by $n$ undershoots the
true variance.
\end{exercise}

\begin{exercise}[Fisher information for the Normal mean]
\label{ch30:exc-fisher-normal}
For $X\sim\Normal(\mu,\sigma^2)$ with $\sigma^2$ known, show the Fisher
information about $\mu$ is $I(\mu)=1/\sigma^2$. Deduce the Cramer--Rao floor for
estimating $\mu$ from $n$ observations, and verify that $\bar X$ attains it.
What does $I(\mu)=1/\sigma^2$ say about how noise affects how much each
observation teaches you?
\end{exercise}

\begin{exercise}[A bias--variance numerical experiment]
\label{ch30:exc-bv}
Consider estimating $p$ from $n=5$ Bernoulli tosses using either
$\hat p = k/5$ or the shrunk $\tilde p=(k+1)/7$. For the true value $p=0.2$:
(a) list the six possible values of $k$ and their binomial probabilities;
(b) compute $\E[\hat p]$, $\E[\tilde p]$ and confirm $\hat p$ is unbiased while
$\tilde p$ is biased toward $1/2$;
(c) compute the MSE of each and report which estimator wins at $p=0.2$.
\end{exercise}

\begin{exercise}[Invariance]
\label{ch30:exc-invariance}
Let $X_i\sim\Expo(\lambda)$ with MLE $\hat\lambda=1/\bar X$. Using the invariance
property, write down the MLE of (a) the mean waiting time $1/\lambda$, (b) the
median waiting time $(\log 2)/\lambda$, and (c) the probability $e^{-\lambda t}$
that the next gap exceeds a fixed time $t$. No new differentiation allowed.
\end{exercise}

\begin{exercise}[Two-parameter score, conceptual]
\label{ch30:exc-two-param}
For $\Normal(\mu,\sigma^2)$ with both parameters unknown, you derived
$\hat\mu=\bar X$ and $\hat\sigma^2=\frac1n\sum(X_i-\bar X)^2$. Explain why the
estimate of $\mu$ did not depend on $\sigma^2$ but the estimate of $\sigma^2$
depended on $\mu$ (through $\bar X$). What does this asymmetry say about the
shape of the likelihood surface, and which parameter is ``easier'' to pin down?
\end{exercise}

\begin{exercise}[Consistency from MSE -- a proof]
\label{ch30:exc-consistency}
Prove that if $\operatorname{Bias}(\hat\theta_n)\to 0$ and
$\Var(\hat\theta_n)\to 0$ then $\hat\theta_n \pto \theta$. (Hint: write
$\mathrm{MSE}=\Var+\operatorname{Bias}^2$ and apply Chebyshev's / Markov's
inequality to $(\hat\theta_n-\theta)^2$.) Then give one example of a consistent
estimator that is biased at every finite $n$.
\end{exercise}

\begin{exercise}[Why the likelihood is not a probability over $\theta$]
\label{ch30:exc-not-prob}
For $\Bern(p)$ with data ``$1$ head in $1$ toss,'' the likelihood is $L(p)=p$.
(a) Compute $\int_0^1 L(p)\,\dd p$ and confirm it is not $1$.
(b) Explain in your own words why this does not violate any axiom of probability,
and what would have to change (philosophically and mathematically) for an
expression like ``the probability that $p$ lies in $[0.4,0.6]$'' to make sense.
(This is the doorway to the next chapter.)
\end{exercise}

\begin{exercise}[Maximum likelihood as closeness to the truth]
\label{ch30:exc-kl}
Argue informally that maximizing the average log-likelihood
$\frac1n\sum_i \log f(X_i\given\theta)$ is, for large $n$, the same as minimizing
the Kullback--Leibler divergence $\KL(f_{\theta_0}\,\|\,f_\theta)$ from the true
distribution to the candidate. (Hint: by the law of large numbers
$\frac1n\sum_i\log f(X_i\given\theta)\to \E_{\theta_0}[\log f(X\given\theta)]$;
subtract the $\theta$-free quantity $\E_{\theta_0}[\log f(X\given\theta_0)]$ and
recognize what remains.) What does this say about what the MLE is ``really''
doing?
\end{exercise}
